\documentclass[../../main.tex]{subfiles}% 注意这里的文件路径不能用 ./main.tex ,否则用latexmk编译子文件会报错
\graphicspath{{\subfix{./image/}}{\subfix{../../image/}}} % 指定图片引用路径,后续可以直接使用图片文件名
% 如果图片存放在上级目录,则使用 ../../image/ 作为备用路径

% 例如:
% \begin{figure}[H]
% \centering
% \includegraphics[width=0.3\textwidth]{图.png}
% \caption{}
% \label{figure:图}
% \end{figure}
% 注意:上述\label{}一定要放在\caption{}之后,否则引用图片序号会只会显示??.

\begin{document}

% SOURCE p=31; printed=16; object=ch00-sec07-title; region=bottom
\section{注记与参考文献}\label{section:ch00-notes-references}

% SOURCE p=31; printed=16; object=ch00-notes-set-theory-start; region=bottom
\cref{section:ch00-set-language,section:ch00-orderings,section:ch00-cardinality,section:ch00-well-ordered-sets}：对初学者而言，集合论最好的讲解是 Halmos\cite{folland-bib-062}，而 Smullyan 和 Fitting\cite{folland-bib-135} 是较高层次的一本好教材。Kelley\cite{folland-bib-083}
% SOURCE p=32; printed=17; object=ch00-notes-set-theory-end; region=top
也简明地介绍了公理集合论的基础。上述各书都从选择公理推导了 Hausdorff 极大原理，Hewitt 和 Stromberg\cite{folland-bib-076} 也作了这一推导。

% SOURCE p=32; printed=17; object=ch00-notes-choice; region=upper
选择公理（或与之等价的某个命题）通常被视为集合论公理化表述中的基本公设之一。一些持直觉主义或构造主义观点的数学家拒绝接受它，其理由是：除非以某种相当明确的方式说明如何构造一个数学对象，否则就不能说已经证明该对象存在；而选择公理的全部意义，正是在构造方法失效（或过于繁琐）时提供存在性定理。真正受这些异议困扰的人属于少数，作者本人不在其中；本书使用选择公理虽不频繁，但也不加限制。

% SOURCE p=32; printed=17; object=ch00-notes-continuum-hypothesis; region=middle
\textbf{连续统假设}断言：若 $\card(X)<\mathfrak{c}$，则 $X$ 可数。（由可数序数集 $\Omega$ 的构造容易得到，对任意不可数集 $X$，都有 $\card(\Omega)\leqslant\card(X)$；因此，等价的断言是 $\card(\Omega)=\mathfrak{c}$。）由 Gödel 和 Cohen 的工作可知，假定集合论的标准公理（包括选择公理）本身相容，则连续统假设及其否定都与这些公理相容。（关于选择公理和连续统假设的相容性与独立性定理，可参见 Smullyan 和 Fitting\cite{folland-bib-135} 的讲解。）有些数学家愿意接受连续统假设为真，似乎是为了方便；但 Gödel\cite{folland-bib-056} 和 Cohen\cite[151]{folland-bib-026} 都曾怀疑它应当为假，而截至本书写作时，尚无人为任何一方找到真正有说服力的证据。我个人的看法是：若一个问题的答案竟然依赖于连续统假设，就应放弃这个问题，转而提出另一个问题；如果集合论出现重大突破，这一看法也可能改变。

% SOURCE p=32; printed=17; object=ch00-notes-metric-spaces; region=lower
\cref{section:ch00-metric-spaces}：关于度量空间更详细的讨论，可参见 Loomis 和 Sternberg\cite{folland-bib-095} 以及 DePree 和 Swartz\cite{folland-bib-032}。
\end{document}
